\documentclass{Article}

\usepackage[utf8]{inputenc}
\usepackage[T1]{fontenc}
\usepackage[francais]{babel}
\usepackage{xcolor}
\usepackage[top=2cm,left=2cm,right=2cm,bottom=2cm]{geometry}
\usepackage[pdfborder={0 0 0}]{hyperref}
\usepackage{graphicx}

\begin{document}

\begin{titlepage}

	\includegraphics[scale=0.4]{logo-ara.png}
	\hfill
	\includegraphics[scale=0.15]{logo-exia.png}

	\vfill
	\centering
	{\huge Mise en place d'un environnement Docker et \\ développement d'une solution e-commerce}\par
	\vspace{1cm}
	\begin{minipage}{0.45\textwidth}
	\raggedleft
	{\Large Rapport de stage de 2\up{e} année}\par
	\vspace{0.25cm}
	{\Large ARA }\par
	\vspace{0.25cm}
	{\Large Administrateurs Réseaux Associés}\par
	\end{minipage}\hspace{1.5cm}%
	\begin{minipage}{0.2\textwidth}
	\raggedright
	\textit{Stagiaire}\par
	Baptiste \bsc{Saclier}\par
	\vspace{0.25cm}
	\textit{Maître de stage}\par
	Alexandre \bsc{Pautet}\par
	\vspace{0.25cm}
	\textit{Tuteur de formation}\par
	Thierry \bsc{Blanc}\par
	\end{minipage}
	\vfill

	\raggedleft
	{\large Le \today}
\end{titlepage}

\clearpage
\newpage\null\thispagestyle{empty}\newpage

\tableofcontents

\section{Remerciements}
\section{Présentation de l'entreprise}
	\subsection{Clients}
	\subsection{Organigramme}
\section{Présentation de l'environnement de travail}

	Pour la création de l'outil demandé, l'entreprise m'a fourni un ordinateur spécifique pré-configuré
dédié à mon travail.

	\subsection{Environnement de développement}

	Pour le développement de l'outil, aucun IDE\footnote{Integrated development environment : Environnement tout intégré permettant la conception,la documentation et la compilation d'une application} n'était imposé, je me suis donc dirigé vers celui que je connaissais et que j'ai utilisé durant ma formation soit PHP Storm sur Microsoft Windows. De plus pour m'assister dans la construction et la gestion des versions j'ai utilisé SVN.

	\subsection{Cycle de développement}

	%TODO to do

\section{Projet Magento Docker}

	Le projet Magento Docker fut ma première mission au sein de l'entreprise. L'objectif de cette mission est de mettre à disposition des développeurs une installation permettant de tester le code sans dépendance a un serveur unique.

	\subsection{Problématiques}

	Avant la mise en place de l'environnement docker par le biais de cette mission les développeurs utilisant les technologies du projet Aello ne pouvaient pas tester sur leur machine due a une incompatibilité entre Magento et Windows. Ainsi, tout les tests étaient effectué sur une machine dite "Sandbox" commune à tout les membres de l'équipe. Cela ne pose pas de gros problèmes avec une équipe de 2 personnes mais le temps de développement peut être allongé inutilement avec une équipe plus importante. En effet, si les changements d'un développeur rendent l'application "Sandbox" instable, tout les développeurs qui testent leur code sur cette application au même moment sont impactés par cette instabilité et donc bloqués jusqu’à la correction de l’erreur.

	\subsection{Image mise en place}
	\subsection{Scripts de gestion}
		\subsubsection{NinjaSync}
	\subsection{Conclusion}
\section{Projet Aello}
	\subsection{Technologies utilisés}
		\subsubsection{Magento}
		\subsubsection{Docker}
	\subsection{Projet initial}
		\subsubsection{Présentation du client}
		\subsubsection{Architecture}
	\subsection{Développement du lot 2}
		\subsubsection{Corrections et amélioration visuelle}
		\subsubsection{Produits favoris}
		\subsubsection{Produits récurrents}
		\subsubsection{Synchronisation des factures}
		\subsubsection{Synchronisation ERP}
	\subsection{Conclusion}	
\section{Bilan de stage}
	\subsection{Difficultés rencontrés}
	\subsection{Bilan personnel}

\end{document}